\section{Speech Translation Architectures and System Comparisons}
\label{sec:architectures}

\subsection{Architectures: Direct vs. Cascaded Paradigms}

\begin{figure}[!ht]
\centering
\resizebox{\columnwidth}{!}{%
\begin{tikzpicture}[
  node distance=6mm and 6mm,
  data/.style={draw, rounded corners=2pt, minimum height=7mm,
    minimum width=17mm, align=center, font=\small, inner sep=2pt},
  proc/.style={draw, minimum height=7mm, minimum width=13mm,
    align=center, font=\small, inner sep=2pt},
  lbl/.style={font=\small\bfseries},
  arrow/.style={-{Stealth[length=2mm]}, thick}
]

% Cascaded pipeline
\node[data] (audio1) {Speech};
\node[proc, right=of audio1] (asr) {ASR};
\node[data, right=of asr] (text) {Transcript};
\node[proc, right=of text] (mt) {MT};
\node[data, right=of mt] (trans1) {Translation};

\draw[arrow] (audio1) -- (asr);
\draw[arrow] (asr) -- (text);
\draw[arrow] (text) -- (mt);
\draw[arrow] (mt) -- (trans1);

\node[lbl, above=1mm of text] {Cascaded ST};

% Direct pipeline
\node[data, below=11mm of audio1] (audio2) {Speech};
\node[proc, right=of audio2, minimum width=56mm] (direct)
  {Direct ST};
\node[data, right=of direct] (trans2) {Translation};

\draw[arrow] (audio2) -- (direct);
\draw[arrow] (direct) -- (trans2);

\node[lbl, above=1mm of direct] {Direct ST};

\end{tikzpicture}%
}
\caption{Simplified inference pipelines based on Sperber and
Paulik~\cite{sperber2020endtoend}. The cascade passes a single best
source-language transcript from ASR to MT. Direct ST produces target-language
text without an intermediate transcript.}
\label{fig:cascade-direct}
\end{figure}

Cascaded speech translation (ST) works as a two-step pipeline. An automatic
speech recognition (ASR) module first transcribes the audio into text in the
source language, and a machine translation (MT) model then translates that
transcript into the target language. The advantage of this separation is
practical. The two components can be developed, trained and replaced
independently. Each can also use the corpora already collected for its own
task: ASR draws on speech paired with transcripts, MT on parallel text, and
both are more widely available than audio paired with
translations~\cite{sperber2020endtoend}. Fig.~\ref{fig:cascade-direct}
contrasts this pipeline with the direct alternative described below.

The same separation creates the weakness most often associated with this
design: error propagation. Where the interface between the two components is
a single best transcript, recognition errors can pass into the MT stage, and
the MT model cannot return to the acoustic signal to check what was said.
The size of this effect varies. Le et al.~\cite{le2017disentangling} studied
French-to-English speech translation and found that improving the ASR
component from 21.86\% to 16.90\% word error rate on their development set
raised BLEU from 26.73 to 28.89. A few recognition errors sometimes caused
many translation errors, whereas a larger number, including morphological
substitutions, could have little effect. Word error rate alone therefore
does not show how much recognition errors affect translation quality,
although their pipeline used phrase-based statistical MT and these effects
need not carry over to neural systems.

Error propagation can be reduced. Sperber and
Paulik~\cite{sperber2020endtoend} list lattices, confusion networks and
error-tolerant MT as established countermeasures, and the cascade of Bentivogli
et al.~\cite{bentivogli2021cascade} fine-tunes its MT component on a mixture of
human and automatic transcripts for the same reason. Two other problems arise
at the interface between the components. MT models trained mainly on written
text may struggle with the disfluent, spoken-style output that ASR produces, a
mismatch between MT training text and ASR output. And once the pipeline commits
to a transcript, anything the transcript does not encode, including prosody, is
lost to the MT model~\cite{sperber2020endtoend}.

Direct ST architectures bypass the intermediate transcript. Weiss et
al.~\cite{weiss2017sequence} built an early direct sequence-to-sequence
system of this kind on 163 hours of Spanish telephone speech paired with
English translations, mapping acoustic features onto target-language text in
a single network. Removing the transcript interface removes the error
propagation that comes with it, and a single model can in principle run with
lower latency than two chained models, although Weiss et al. did not report
latency measurements. Paired speech and translation data is scarce, and
harder to obtain than the separate ASR and MT corpora available to cascaded
systems. Sperber and Paulik~\cite{sperber2020endtoend} describe this as a
possible trade-off between data efficiency and modelling power, and note
that the techniques direct models use to compensate, such as training on ASR
and MT data, may reintroduce some of the issues associated with cascades.
They present this as a question for further analysis rather than a settled
result.

\subsection{Empirical System Comparisons}

Bentivogli et al.~\cite{bentivogli2021cascade} compared cascaded and direct
systems built on the same core Transformer technology, drawing on overlapping
training resources used through different pipelines, and evaluated both on the
same MuST-C test segments for English to German, Spanish and Italian. They used
automatic metrics together with professional post-editing, assessed using two
measures based on translation edit rate (TER): human-targeted TER (HTER) and
multi-reference TER (mTER). On those post-editing measures the two approaches
performed similarly for English-to-German and English-to-Italian. For
English-to-Spanish the direct system was significantly better, while BLEU
computed against independent references favoured the cascade. The authors note
that these measures can lead to different conclusions.

A more detailed analysis showed different error patterns across directions. The
cascade made fewer morphological errors in English-to-German and
English-to-Italian, and fewer word-order errors in English-to-German. For
English-to-Spanish, the direct system made fewer lexical, morphological and
word-order errors. On a filtered subset prepared for manual annotation, direct
systems left more source words untranslated, though some omissions were
discourse markers that the authors say need not count as errors. The direct
system appeared to handle prosody better and made fewer errors in understanding
speech. The authors call both findings tentative, since few prosody cases could
be isolated and errors inside direct models are harder to observe.

The direct model used an encoder initialised from an ASR model, knowledge
distillation from a teacher MT model, and synthetic data produced by translating
transcripts from ASR training corpora~\cite{bentivogli2021cascade}. The
experiment therefore does not show that a direct model could match the cascade
without ASR and MT resources, nor that the gains it does achieve follow from
the inference architecture alone.

A study of Basque and Spanish reported different results. Etchegoyhen et
al.~\cite{etchegoyhen2022casestudy} compared both directions for this pair,
where Basque has rich morphology and relatively free word order. On their
main test sets the cascade was stronger for Basque-to-Spanish, while a
wav2vec-based direct model scored higher for Spanish-to-Basque. In a manual
ranking on a separate sample of parliamentary speech, evaluators preferred
the cascaded output in every configuration except Spanish-to-Basque with
additional training data, where the two were judged roughly equal. When
results were grouped by ASR word error rate, direct models did not
consistently perform better on subsets with higher WER. The pattern varied
with language direction and training data.

Shared tasks compare more systems, which differ in many respects besides
architecture. In the IWSLT 2021 offline task, cascaded systems took the top two
places, and the best cascade scored 2 BLEU points above the best direct
submission. With the original subtitle-style TED references, a different direct
submission ranked highest among direct systems, and the gap to the best cascade
was 0.7 BLEU. This shows that reference choice can change both system rankings
and the apparent gap between the two approaches. In the multilingual task of
the same campaign, the three teams submitting both architectures found their
direct systems stronger.
The organisers suggest that differences between the tasks in auxiliary data
volume and in audio segmentation may help explain
this~\cite{anastasopoulos2021iwslt}. Two years later, cascaded systems were
again generally stronger in the IWSLT 2023 offline
task~\cite{agarwal2023iwslt}.

More recently, Papi et al.~\cite{papi2026hearing} benchmarked twenty-two
systems, four direct and twelve cascaded alongside six speech-enabled language
models, across sixteen benchmarks and a range of acoustic and linguistic
conditions. Cascaded systems were the most reliable overall, and performance
depended on the condition tested rather than on the architecture alone. In a
human evaluation comparing one model from each group, the direct model produced
more undertranslation than the cascade, consistent with the pattern reported by
Bentivogli et al.~\cite{bentivogli2021cascade}.

\subsection{From System Comparisons to Accent}

Across these studies, no approach performs best in every setting, and the
reported differences vary with the language pair, the evaluation data, the
reference translations and the metric, a dependence examined further in
Section~\ref{sec:evaluation}. Earlier comparisons such as Bentivogli et
al.~\cite{bentivogli2021cascade} and Etchegoyhen et
al.~\cite{etchegoyhen2022casestudy} did not report results by accent group.
Other speech conditions were also considered: Weiss et
al.~\cite{weiss2017sequence} used telephone speech, and Bentivogli et
al.~\cite{bentivogli2021cascade} analysed audio duration and speech rate and
examined errors involving poor audio quality or overlapping speech.

Recent studies report accent-level results. The IWSLT 2024 offline task
added an accent challenge built from the Edinburgh International Accents of
English Corpus, and all submissions to that task used a cascaded
architecture~\cite{ahmad2024iwslt}. Papi et al.~\cite{papi2026hearing}
evaluate accent conditions across direct, cascaded and speech-enabled
systems and suggest that the choice of speech encoder matters for accent
robustness. Section~\ref{sec:accent-st} examines accent robustness in speech
translation in more detail. Comparisons at the system level can show how
translation quality varies across accent groups, but they do not on their
own explain where those differences come from. The next section reviews how
accent affects ASR, which may help explain errors in cascaded speech
translation.
